% 자동 생성 — scripts/build_tables.py. 손으로 고치지 않는다.
% 출처: output/results/06d_pretrend.json : A_기준분기
\begin{table}[htbp]
  \centering
  \caption{기준 분기를 옮겨 본 거래량 효과. $k=-2$ 만 크게 튄다.}
  \label{tab:pretrend-ref}
  \begin{tabular}{lrrrr}
    \toprule
    \textbf{기준 분기} & \textbf{$k{=}0$} & \textbf{$k{=}1$} & \textbf{$k{=}4$} & \textbf{사전 유의 칸} \\
    \midrule
    -1 & -79.8 & -66.8 & -39.7 & 3 \\
    -2 & -61.2 & -36.2 & 16.0 & 6 \\
    -3 & -73.3 & -56.2 & -20.3 & 3 \\
    -4 & -77.8 & -63.4 & -33.5 & 2 \\
    -5 & -78.4 & -64.5 & -35.5 & 2 \\
    -6 & -78.6 & -64.8 & -36.0 & 1 \\
    \bottomrule
  \end{tabular}

  \begin{minipage}{\linewidth}\footnotesize
  \textit{주.} $k=-2$ 는 12·16 대책 직후라 눌려 있는 분기다. 나머지는 $-73 \sim -80\%$ 로 좁게 모인다.
  \end{minipage}
\end{table}
